\documentclass[conference]{IEEEtran}

\usepackage{cite}
\usepackage{amsmath,amssymb,amsfonts}
\usepackage{graphicx}
\usepackage{textcomp}
\usepackage{xcolor}
\usepackage{booktabs}
\usepackage{url}

\begin{document}

\title{Project Title}

\author{
\IEEEauthorblockN{Student Name, Student Name, Student Name}
\IEEEauthorblockA{
\textit{*Department of Advanced Computing Sciences*} \\\\
\textit{*Faculty of Science and Engineering*} \\\\
\textit{*Maastricht University*} \\\\
Maastricht, The Netherlands}
}

\maketitle

\begin{abstract}
% Briefly summarize the complete project:
% - What problem did you address?
% - What approach did you use?
% - What are the most important results?
% - What is the main conclusion?
%
% The abstract should be understandable without reading the full report.
\end{abstract}

\begin{IEEEkeywords}
% Add 3--5 keywords that describe the main topic, methods, or application domain.
\end{IEEEkeywords}

\section{Introduction}
% Introduce the problem and its context.
% Explain why the problem is scientifically, societally, or practically relevant.
% Describe the relevant stakeholder perspective(s).
% State the research objective(s) and research question(s).
% Clarify the scope of the project where necessary.
%
% Use scientific literature to support claims about the problem and its relevance.

\section{Related Work}
% Discuss the scientific literature most relevant to your project.
% Explain what is already known and how existing approaches relate to your work.
% Position your project with respect to the literature.
%
% Do not present papers as an unrelated list. Organize the discussion around
% themes, approaches, findings, or limitations that are relevant to your project.
% This section can also be merged into the Introduction

\section{Methodology}
% Describe the methodology used to answer the research question(s).
% Explain and justify important methodological choices using scientific literature
% where appropriate.
%
% Depending on the project, this section may include:
% - data sources and preprocessing;
% - algorithms, models, or system design;
% - experimental setup;
% - baselines or comparison methods;
% - evaluation protocol;
% - evaluation metrics and success criteria;
% - implementation decisions relevant to the methodology.
%
% Provide enough detail for the reader to understand and assess the approach.

\section{Results}
% Present the results of the project clearly and objectively.
% Structure the section around the research question(s) or evaluation criteria
% where appropriate.
%
% Use tables, figures, charts, or other visualizations when they improve clarity.
% Refer to every figure and table in the text.
%
% Focus here on what was observed. Interpretation belongs mainly in the Discussion.

\section{Discussion}
% Interpret the results in relation to the research question(s) and objective(s).
% Explain what the findings mean.
% Compare and position the results with respect to relevant scientific literature.
% Discuss the validity, reliability, and limitations of the evaluation.
% Discuss unexpected findings where relevant.
%
% Clearly distinguish conclusions supported by evidence from speculation.

\section{Conclusion}
% Summarize the main findings.
% Explicitly answer the research question(s).
% State the main evidence-based conclusions.
% Provide recommendations or directions for future work where appropriate.
%
% Do not introduce substantial new results in this section.

% Optional sections may be added where appropriate for the project, for example:
% \section{Ethical, Legal, Privacy, or Security Considerations}
% \section{System Architecture}
% \section{Implementation}
% \section{Experiments}
%
% Only add sections when they contribute meaningfully to the scientific report.

\bibliographystyle{IEEEtran}
\bibliography{references}

\end{document}
